\documentclass[../../main]{subfiles}

\ifSubfilesClassLoaded{%
    \zexternaldocument*{../1.Category-Theory/exported/aux/main}%
    \zexternaldocument*{../2.Homological-Algebra/exported/aux/main}%
}{}

\begin{document}
\chapter{Presheaves and sheaves}\label{ch:presheaves-and-sheaves}

In previous chapters, we developed the framework of category theory and homological algebra to study mathematical objects and their global transformations.
However, many geometric and topological problems are inherently local.
To bridge the gap between local properties and global structures, we introduce the concepts of presheaves and sheaves.

Intuitively, a presheaf allows us to attach data (such as modules, groups, or rings) to the open sets of a topological space in a way that respects the restriction from larger open sets to smaller ones.
A sheaf goes a step further by requiring that this data satisfies a local-to-global principle: sections that are locally compatible can be uniquely glued together into a global section.

This chapter begins with a motivating example centered around smooth functions, tracking how these concrete structures naturally abstract into general definitions.
We will explore the local properties of sheaves through the construction of stalks, establish the sheafification functor as a left adjoint to the forgetful functor, and conclude by examining the rich homological properties of sheaves valued in an abelian category.

\section{A motivating example}\label{sec:motivating-example-sheaves}
We begin this chapter by giving a concrete example of sheaves following \cite{vakilRISINGSEAFoundations2025} that will help us motivate the abstractions.

Consider smooth (\(\mathcal{C}^{\infty}\)) functions on the topological space \(X = \mathbb{R}^{n}\) (or more generally on a manifold \(X\)).
The sheaf of smooth functions on \(X\) is the data of all smooth functions on all open subsets of \(X\).
Note that for each open subset \(U \subseteq X\) we have an \(\mathbb{R}\)-module of smooth functions \(\sheaf{O}(U)\).

Given an smooth function \(f\) in an open set \(U\) we can restrict it to an open subset \(V \subseteq U\).
In other words, we have a restriction map \(\res_{U, V}\colon \sheaf{O}(U) \to \sheaf{O}(V)\).
Furthermore if we have \(W \subseteq V \subseteq U\) open subsets restricting \(f\) to \(V\) and then to \(W\) is the same as restrinting \(f\) directly to \(W\), which means that the following diagram is commutative
\begin{center}
  \begin{tikzcd}
    \sheaf{O}(U) \arrow[rr, "\res_{U, W}"] \arrow[dr, "\res_{U, V}"'] & & \sheaf{O}(W) \\
    & \sheaf{O}(V) \arrow[ur, "\res_{V, W}"']
  \end{tikzcd}
\end{center}

Notice how if we replace smooth functions with, say, continuous, real-analytic or bounded functions this example would work analogously.
Therefore all of this ``nice'' types of functions share some common properties.
The abstraction of this data --- the collection of \(\mathbb{R}\)-modules plus the restriction maps --- is what we will call a presheaf (of \(\mathbb{R}\)-modules) over \(X\).

However in this particular our presheaf is more rich.
Indeed, lets consider an open cover \(\mathcal{U} \coloneq \{U_{i}\}_{i \in \mathcal{I}}\) of an open set \(U \subseteq X\).

On the one hand if we have two smooth function \(f_{1}, f_{2} \in \sheaf{O}(U)\) such that both agree on every \(U_{i}\) then they must be the same function to begin with.
In the language of the restriction maps we have that if \(\res_{U_{i}, U_{j}} f_{1} = \res_{U_{i}, U_{j}} f_{2}\) for all \(U_{i}, U_{j} \in \mathcal{U}\) then \(f_{1} = f_{2}\).
This means that we can identify functions on an open set by looking at them on a covering by smaller open set.

On the other suppose now we have a collection of smooth function \(f_{i} \in \sheaf{O}(U_{i})\) for all \(i \in \mathcal{I}\) such that they agree on all pairwise overlaps.
Then we can \textit{glue} them together forming an smooth function \(f\) in all of \(U\).
Again, with the notation of the restriction maps we have that if \(\res_{U_{i} \cap U_{j}} f_{i} = \res_{U_{i} \cap U_{j}} f_{j}\) for all \(i, j \in \mathcal{I}\) there exists \(f \in \sheaf{O}(U)\) such that \(\res_{U_{i}} f = f_{i}\).

This two properties are what will differentiate a presheaf from a sheaf.
Note that they still hold for continous or real-analytic functions, but they no longer work in the bounded functions case (the second property fails).

\section{Presheaves and sheaves}

We are now ready to start abstracting this data into a more general setting, where we can use the full power of the category theory we have developed previously.

\begin{definition}[Presheaf]\label{def:presheaf}
  Let \(X\) be a topological space and let \(\category{C}\) be a \hyperref[def:category]{category}.
  The \hyperref[def:category]{category} of \(\category{C}\)-valued \hyperref[def:presheaf]{presheaves} over \(X\) is the \hyperref[def:functor-category]{functor category}
  \begin{equation*}
    \PSh_{X}(\category{C}) \coloneq \Fct{\Open(X)^{\op}}{\category{C}}
  \end{equation*}

  Less abstractly, a \(\category{C}\)-valued \hyperref[def:presheaf]{presheaf} \(\mathcal{F}\) is the collection of data:
  \begin{itemize}
    \item For each open set \(U \subseteq X\) an object \(\mathcal{F}(U) \in \category{C}\) called \textit{sections of \(\mathcal{F}\) over \(U\)}.
    \item For every inclusion of open sets \(V \subseteq U\) a morphism \(\res_{U, V}\colon \mathcal{F}(U) \to \mathcal{F}(V)\) in \(\category{C}\) called \textit{restriction} such that
          \begin{itemize}
            \item For every open set \(U \subseteq X\), \(\res_{U, U} = \id_{U}\).
            \item For every inclusion of open sets \(W \subseteq V \subseteq U\), \(\res_{U, V} \circ \res_{V, W} = \res_{U, W}\).
          \end{itemize}
  \end{itemize}
  A morphism of \hyperref[def:presheaf]{presheaves} \(\phi: \mathcal{F} \to \mathcal{G}\) is a collection of morphisms
  \begin{equation*}
    \{\phi(U)\colon \mathcal{F}(U) \to \mathcal{G}(U)\}_{U \in \Open(X)}
  \end{equation*}
  in \(\category{C}\) making the following diagram commute for all \(V \subseteq U\)
  \begin{center}
    \begin{tikzcd}
      \mathcal{F}(U) \arrow[r, "\phi(U)"] \arrow[d, "\res_{U, V}"'] & \mathcal{G}(U) \arrow[d, "\res_{U, V}"] \\
      \mathcal{F}(V) \arrow[r, "\phi(V)"'] & \mathcal{G}(V)
    \end{tikzcd}
  \end{center}
\end{definition}

\begin{example}
  Smooth, real-analytic and bounded functions on \(X = \mathbb{R}^{n}\) (or more generally in a manifold \(X\)) are \hyperref[def:presheaf]{presheaves}.
\end{example}


Since \hyperref[def:presheaf]{presheaves} are particular instances of \hyperref[def:functor]{functors} we have already seen several examples back in \autoref{ch:functors-natural-transformations}.
One example that will be useful later is the following.
\begin{definition}[Constant presheaf]\label{def:constant-presheaf}
  Let \(X\) be a topological space.
  For any object \(A \in \category{C}\) the \textbf{constant presheaf} with value \(A\), denoted \(\underline{A}_{X}^{\operatorname{pre}}\) is the \hyperref[def:presheaf]{\(\category{C}\)-valued presheaf} associated to the \hyperref[ex:constant-functor]{constant functor}
  \begin{center}
    \begin{tikzcd}
      \llap{\(\underline{A}_{X}^{\operatorname{pre}}\colon\)} \Open(X)^{\op} \arrow[r] & \category{C} \\[-2em]
      U \arrow[d, "V \subseteq U"', ""{name=source, right}] & A \arrow[d, "\id_{A}" ""{name=target, left}] \arrow[from=source, to=target, mapsto] \\
      V  & A
    \end{tikzcd}
  \end{center}

  Note that this defines a \hyperref[def:functor]{functor}
  \begin{align*}
    \underline{-}_{X}^{\operatorname{pre}}\colon \category{C} \to \PSh_{X}(\category{C})
  \end{align*}
\end{definition}

We now turn our attention to a particular class of \hyperref[def:presheaf]{presheaves}.

Instead of working on a general \hyperref[def:category]{category} \(\category{C}\), to simplify the exposition we will work with \hyperref[def:category]{categories} which are ``sets with an algebraic structure'' (for example \(\Grp, \Ab, \Ring, \ldots\) but not \(\Top\)).

\begin{definition}\label{def:algebraic-structure}
  A \textbf{type of algebraic structure} is a \hyperref[def:category]{category} \(\category{C}\) together with a \hyperref[def:functor]{functor} \(F\colon \category{C} \to \Set\) verifying:
  \begin{enumerate}
    \item \(F\) is \hyperref[def:functor-types]{faithful}.
    \item \(\category{C}\) has \hyperref[def:categorical-limit]{limits} and \(F\) commutes with them.
    \item \(\category{C}\) has \hyperref[def:direct-limit]{direct limits} and \(F\) commutes with them.
    \item \(F\) reflects \hyperref[def:categorical-isomorphism]{isomorphisms}.
  \end{enumerate}
\end{definition}

\begin{example}
  The \hyperref[def:category]{categories} of pointed sets, (abelian) groups, monoids, rings \(R\)-modules, and Lie algebras -- together with their respective forgetful \hyperref[def:functor]{functors} -- are \hyperref[def:algebraic-structure]{types of algebraic structures}.
\end{example}

We state some properties of this type of \hyperref[def:category]{categories} without proof.
\begin{lemma}[{\cite[\href{https://stacks.math.columbia.edu/tag/007O}{Tag 007O}]{stacks-project}}]\label{lemma:algebraic-structure-properties}
  Let \((\category{C}, F)\) be a \hyperref[def:algebraic-structure]{type of algebraic structure}. Then:
  \begin{enumerate}
    \item \(\category{C}\) has a \hyperref[def:initial-terminal-object]{terminal object} \(1\) and \(F(1) = \{\ast\}\).
    \item \(\category{C}\) has \hyperref[def:categorical-product]{products} and \(F \left( \prod X_{i} \right) = \prod F \left( X_{i} \right)\).
    \item \(\category{C}\) has \hyperref[def:pullback]{pullbacks} and \(F(X \times_{Z} Y) = F(X) \times_{F(Z)} F(Y)\).
    \item \(\category{C}\) has \hyperref[def:equaliser]{equalisers} and if \(\operatorname{eq} \to X\) is the \hyperref[def:equaliser]{equaliser} of \(f, g\colon X \rightrightarrows Y\) then \(F(\operatorname{eq}) \to F(X)\) is the \hyperref[def:equaliser]{equaliser} of \(Ff, Fg\colon F(X) \rightrightarrows F(Y)\).
    \item \(X \to Y\) is a \hyperref[def:mono-epi-morphism]{monomorphism} if and only if \(F(X) \to F(Y)\) is injective.
    \item If \(Ff\colon F(X) \to F(Y)\) is surjective then \(f\) is an \hyperref[def:mono-epi-morphism]{epimorphism}.
    \item Given \(X_{1} \to X_{2} \to \cdots\), \(\varinjlim X_{i}\) exists and \(F(\varinjlim X_{i}) = \varinjlim F(X_{i})\). And more generally for any \hyperref[def:direct-limit]{direct limit}.
    \item \hyperref[def:direct-limit]{Direct limits} commute with \hyperref[def:categorical-limit]{limits}.
  \end{enumerate}
\end{lemma}

Therefore in everything that follows after this point we will implicitly assume that \(\category{C}\) is a \hyperref[def:algebraic-structure]{type of algebraic structure} every time we talk about sheaves.

\begin{definition} \label{def:sheaf}
  Let \(X\) be a topological space.
  The \hyperref[def:category]{category} of \textbf{\(\category{C}\)-valued sheaves} over \(X\) is the \hyperref[def:subcategory]{full subcategory} \(\Sh_{X}(\category{C})\) of \(\PSh_{X}(\category{C})\) whose objects verify, for every open set \(U \subseteq X\) and every open cover \(\{U_{i}\}_{i \in \mathcal{I}}\) of \(U\)
  \begin{enumerate}[label=SH.\arabic*), ref=SH.\arabic*]
    \item\label{axiom:sheaf-identity} \textbf{Identity}: Given two sections \(s, t \in \sheaf{F}(U)\) such that for every \(i \in \mathcal{I} \quad s|_{U_{i}} = t|_{U_{i}}\) we have that \(s = t\).
    \item\label{axiom:sheaf-gluability} \textbf{Gluability}: Given a collection of sections \(\{s_{i} \mid s_{i} \in \sheaf{F}(U_{i})\}_{i \in \mathcal{I}}\) that agree on pairwise intersections (\(\forall i, j \in \mathcal{I} \;\; s_{i}|_{U_{i} \cap U_{j}} = s_{j}|_{U_{i} \cap U_{j}}\)) there exists a section \(s \in \sheaf{F}(U)\) such that \(\forall i \in \mathcal{I} \;\; s|_{U_{i}} = s_{i}\).
  \end{enumerate}
\end{definition}

\begin{notation}
  From now on, to distinguish a \hyperref[def:presheaf]{presheaf} from a \hyperref[def:sheaf]{sheaf} we will denote the former by \(\mathcal{F}, \mathcal{G}, \mathcal{H}, \ldots\) and the latter by \(\sheaf{F}, \sheaf{G}, \sheaf{H}, \ldots\)
\end{notation}

\begin{example}
  Smooth and continuous function form a \hyperref[def:sheaf]{sheaf of \(\mathbb{R}\)-modules} over \(X \coloneq \mathbb{R}\).
  However, bounded functions \(\mathcal{B}\) only form a \hyperref[def:presheaf]{presheaf}.

  Indeed if we consider the open covering \(\mathcal{U} \coloneq \left\{ ]n, n+1[ \right\}_{n \in \mathbb{N}}\) of \(\mathbb{R}\) and \(f_{n}(x) \coloneq n\) we have that \(f_{n} \in \mathcal{B}(]n, n+1[)\) for all \(n \in \mathbb{N}\) but any \textit{glueing} of these functions is clearly not bounded and thus \ref{axiom:sheaf-gluability} fails.
\end{example}

\begin{remark}\label{rem:sheaves-as-equalisers}
  Axioms \ref{axiom:sheaf-identity} and \ref{axiom:sheaf-gluability} can be interpreted as the following \hyperref[def:equaliser]{equaliser diagram}:
  \begin{center}
    \begin{tikzcd}
      \sheaf{F}(U) \arrow[r] & \displaystyle\prod_{i \in \mathcal{I}} \sheaf{F}(U_{i}) \arrow[r, shift left] \arrow[r, shift right] & \displaystyle\prod_{(i, j) \in \mathcal{I} \times \mathcal{I}} \sheaf{F}(U_{i} \cap U_{j})
    \end{tikzcd}
  \end{center}
  where \(\{U_{i}\}_{i \in \mathcal{I}}\) is an open cover of \(U\). 
  
  In particular taking the empty cover of \(\varnothing \in \Open(X)\), since \hyperref[rem:empty-product-is-terminal]{the empty product is terminal} we have
  \begin{center}
    \begin{tikzcd}
      \sheaf{F}(\varnothing) \arrow[r] & 1  \arrow[r, shift left] \arrow[r, shift right] & 1
    \end{tikzcd}
  \end{center}
  i.e: \(\sheaf{F}(\varnothing)\) is a \hyperref[def:initial-terminal-object]{terminal object} of \(\category{C}\).
\end{remark}

\begin{example}[Sheaves on a point]\label{ex:one-point-sheaves-are-category}
  Let \(\{\ast\}\) be the one-point topological space.
  We are going see that there is an \hyperref[def:equivalent-categories]{equivalence of categories}
  \begin{equation*}
    \Sh_{\{\ast\}}(\category{C}) \simeq \category{C}
  \end{equation*}

  Consider the \hyperref[def:functor]{functors}
  \begin{align*}
    F\colon\Sh_{\{\ast\}}(\category{C}) \to \category{C} \qquad G\colon\category{C} \to \Sh_{\{\ast\}}(\category{C})
  \end{align*}
  where \(F(\sheaf{F}) \coloneq \sheaf{F}(\{\ast\})\) for all \(\sheaf{F} \in \Sh_{\{\ast\}}(\category{C})\) and for all \(X \in \category{C}\), \(G(X)\) is the \hyperref[def:sheaf]{sheaf} defined as
  \begin{equation*}
    G(X)(U) \coloneq \begin{cases}
      X & U = \{\ast\} \\
      1 & U = \varnothing
    \end{cases}
  \end{equation*}
  with the obvious restriction maps.
  We have
  \begin{equation*}
    \left( F \circ G \right)(X) = G \qquad \left( G \circ F \right)(\sheaf{F}) \cong \sheaf{F}
  \end{equation*}
  where the \hyperref[def:categorical-isomorphism]{isomorphism} comes from \(\sheaf{F}(\varnothing)\) being \hyperref[def:initial-terminal-object]{terminal} (but not necessarilly \(1\)).
\end{example}

\subsection{Stalks}

The main difference between \hyperref[def:sheaf]{sheaves} and \hyperref[def:presheaf]{presheaves} is their local nature.
We are going to see that there are ``global'' properties of \hyperref[def:sheaf]{sheaves} that can be checked if we know all of the ``local'' information.

Before trying to abstract this, let's go back to our \hyperref[sec:motivating-example-sheaves]{motivating example} of smooth functions on \(\mathbb{R}^{n}\).
A germ of a smooth function at a point \(p \in X\) is an object of the form \((f, U)\) such that \(p \in U\) open and \(f \in \sheaf{O}(U)\), modulo the relation \((f, U) \sim (g, V)\) if there exists \(p \in W \subseteq U \cap V\) open such that \(f|_{W} = g|_{W}\).
We call this set of germs the stalk of \(\sheaf{O}\) at \(p\) and denote it by \(\sheaf{O}_{p}\).
It is easy to see that \(\sheaf{O}_{p}\) has the structure of an \(\mathbb{R}\)-module.
Note that here we have not used the two \hyperref[def:sheaf]{sheaf properties}, so this definition works considering \(\sheaf{O}\) either as a \hyperref[def:presheaf]{presheaf} or as a \hyperref[def:sheaf]{sheaf}.

Furthermore if \(p \in U\) we get maps \(\sheaf{O}(U) \to \sheaf{O}_{p}\), which hints us that we can model this as a \hyperref[def:categorical-colimit]{colimit}.
\begin{definition}\label{def:stalk}
  Let \(X\) be a topological space and let \(\mathcal{F}\) be a \hyperref[def:presheaf]{\(\category{C}\)-valued presheaf}.
  For all points \(p \in X\) we define the \textbf{stalk} of \(\mathcal{F}\) at \(p\) as
  \begin{equation*}
    \mathcal{F}_{p} \coloneq \varinjlim_{\substack{U \ni p \\ U \in \Open(X)}} \mathcal{F}(U)
  \end{equation*}
  Note that this comes with maps \(\mathcal{F}(U) \to \mathcal{F}_{p}\) for all \(U \subseteq X\) open such that \(p \in U\).
  If \(p \in U\) and \(s \in \sheaf{F}(U)\) then the image of \(s\) in \(\sheaf{F}_{p}\) is called the \textbf{germ} of \(s\) at \(p\) and denote it by \(s_{p}\).
  
  Since this is a \hyperref[def:direct-limit]{direct limit} following \autoref{ex:direct-limits-set} we may describe \(\mathcal{F}_{p}\) as
  \begin{equation*}
    \mathcal{F}_{p} = \quot{\{ (U, s) \mid p \in U, s \in \mathcal{F}(U)\}}{\sim}
  \end{equation*}
  where
  \begin{equation*}
    (U, s) \sim (V, t) \iff \exists W \subseteq U \cap V \text{ open} \mid p \in W \land s|_{W} = t|_{W}
  \end{equation*}
  
  Using this characterization we may define the stalk of a \hyperref[def:presheaf]{morphism of presheaves} \(\phi: \mathcal{F} \to \mathcal{G}\) as
  \begin{align*}
    \phi_{x}\colon& \mathcal{F}_{x} \to \mathcal{G}_{x} \\
            & [(U, s)] \mapsto [(U, \phi(s))]
  \end{align*}

  Therefore we have a stalk \hyperref[def:functor]{functor}
  \begin{equation*}
    -_{p}\colon \PSh_{X}(\category{C}) \to \category{C}
  \end{equation*}
  By using the same definition we can extend this to a \hyperref[def:functor]{functor}
  \begin{equation*}
    -_{p}\colon \Sh_{X}(\category{C}) \to \category{C}
  \end{equation*}
\end{definition}

\begin{remark}\label{rem:stalk-commute-finite-limits}
  By the \hyperref[lemma:algebraic-structure-properties]{properties of algebraic structures}, \hyperref[def:stalk]{stalks} commute with \hyperref[def:categorical-limit]{finite limits}.
\end{remark}

The following lemma somewhat explains the agricultural terminology: sections of a \hyperref[def:sheaf]{sheaf} can be seen as a collection of \hyperref[def:stalk]{germs} bundled together.
\begin{lemma}\label{lemma:sheaf-determined-by-germs}
  Let \(\sheaf{F}\) be a \hyperref[def:sheaf]{\(\category{C}\)-valued sheaf} over a topological space \(X\).
  For all open subsets \(U \subseteq X\) the map \(\sheaf{F}(U) \to \prod_{p \in U} \sheaf{F}_{p}\) is injective.

  In other words, a \hyperref[def:sheaf]{sheaf} is completely determined by its germs.
\end{lemma}
\begin{proof}
  Let \(s, t \in \sheaf{F}(U)\) such that for all \(s_{p} = t_{p}\) in \(\sheaf{F}_{p}\) for all \(p \in U\).

  By \hyperref[def:stalk]{definition of stalk} there exists an open neighbourhood \(V_{p} \subseteq U\) of \(p\) such that \(s|_{V_{p}} = t|_{V_{p}}\).
  Since this works for all \(p \in U\), we have an open cover \(\{V_{p}\}_{p \in U}\) of \(U\).

  Therefore since \(\sheaf{F}\) is a \hyperref[def:sheaf]{sheaf} by \ref{axiom:sheaf-identity} we have \(s = t\).
\end{proof}

Using this result we can now prove one of the most important theorems for \hyperref[def:sheaf]{sheaves}.
Roughly speaking, it states that knowing all of the ``local'' information is the same as knowing the ``global'' information.
\begin{theorem}\label{thm:sheaf-morphism-determined-by-stalks}
  Let \(\phi, \psi\colon \mathcal{F} \to \sheaf{G}\) be a \hyperref[def:presheaf]{morphism of presheaves} where \(\mathcal{F}\) is a \hyperref[def:presheaf]{\(\category{C}\)-valued presheaf} and \(\sheaf{G}\) is a \hyperref[def:sheaf]{\(\category{C}\)-valued sheaf}.
  Then \(\phi = \psi\) if and only if for all \(p \in X\) the \hyperref[def:stalk]{stalks at \(p\)} are equal, i.e: \(\phi_{p} = \psi_{p}\).
\end{theorem}
\begin{proof}
  The direction \([\implies]\) is trivial.

  For the \([\impliedby]\) direction suppose that \(\phi_{p} = \psi_{p}\) for all \(p \in X\) and consider the commutative diagram
  \begin{center}
    \begin{tikzcd}
      \mathcal{F}(U) \arrow[r, shift left, "\phi"] \arrow[r, shift right, "\psi"'] \arrow[d, "i_{\mathcal{F}}"'] & \sheaf{G}(U) \arrow[d, tail, "i_{\sheaf{G}}"] \\
      \displaystyle\prod_{p \in U} \mathcal{F}_{p}  \arrow[r, shift left, "\prod \phi_{p}"] \arrow[r, shift right, "\prod \psi_{p}"']& \displaystyle\prod_{p \in U} \sheaf{G}_{p}
    \end{tikzcd}
  \end{center}
  For all \(s \in \mathcal{F}(U)\) we have
  \begin{equation*}
    i_{\sheaf{G}}(\phi(U)(s)) = \left( \phi_{p}(s_{p}) \right)_{p \in U} = \left( \psi_{p}(s_{p}) \right)_{p \in U} = i_{\sheaf{G}}(\psi(U)(s))
  \end{equation*}
  Since \hyperref[lemma:sheaf-determined-by-germs]{\(i_{\sheaf{G}}\) is injective} we have
  \begin{equation*}
    \phi(U)(s) = \psi(U)(s)
  \end{equation*}
  Therefore since this holds for all \(s \in \mathcal{F}(U)\) and \(U \subseteq X\) open we have that \(\phi = \psi\).
\end{proof}

\begin{example}
  \autoref{thm:sheaf-morphism-determined-by-stalks} is false in the case of \hyperref[def:presheaf]{presheaves}.

  Indeed, if we consider \(\mathcal{B}\) the \hyperref[def:presheaf]{presheaf} of bounded functions on \(\mathbb{R}\) and as morphisms the identity \(\id_{\mathcal{B}}\colon \mathcal{B} \to \mathcal{B}\) and \(\psi\colon \mathcal{B} \to \mathcal{B}\) defined as \(\psi(U) = \id_{U}\) if \(U\) is bounded and the zero morphism otherwise since \hyperref[def:stalk]{stalks} are computed by looking at small neighbourhoods and every point \(p \in \mathbb{R}\) admits infinitely small bounded neighbourhoods we have
  \begin{equation*}
    (\id_{\mathcal{B}})_{p} = \psi_{p} = \id_{\mathcal{B}_{p}}
  \end{equation*}
  However, if we consider \(f(x) \coloneq 1 \in \mathcal{B}(\mathbb{R})\) we get
  \begin{equation*}
    \id_{\mathcal{B}}(\mathbb{R})(f) = 1 \neq 0 = \psi(\mathbb{R})(f)
  \end{equation*}
\end{example}

The \hyperref[def:sheaf]{category of sheaves} also has a nice characterization \hyperref[def:categorical-isomorphism]{isomorphism}, \hyperref[def:mono-epi-morphism]{monomorphism} and \hyperref[def:mono-epi-morphism]{epimorphisms} in terms of \hyperref[def:stalk]{stalks} that will be particularly useful once we study \hyperref[def:exact-sequence]{exact sequences} of \hyperref[def:sheaf]{sheaves}.

\begin{proposition}\label{prop:sheaf-isomorphism-determined-at-stalks}
  Let \(\phi\colon \sheaf{F} \to \sheaf{G}\) be a \hyperref[def:sheaf]{morphism of sheaves}.
  Then \(\phi\) is a \hyperref[def:categorical-isomorphism]{isomorphism} (resp. \hyperref[def:mono-epi-morphism]{mono/epimorphism}) if and only if for all \(p \in X\) the \hyperref[def:stalk]{induced stalk map} \(\phi_{p}\colon \sheaf{F}_{p} \to \sheaf{G}_{p}\) is an \hyperref[def:categorical-isomorphism]{isomorphism} (resp. \hyperref[def:mono-epi-morphism]{mono/epimorphism}).
\end{proposition}
\begin{proof}
  We will check the \hyperref[def:categorical-isomorphism]{isomorphism} case.
  \begin{itemize}
    \item[\(\implies\)] Suppose that \(\phi: \sheaf{F} \to \sheaf{G}\) is an \hyperref[def:categorical-isomorphism]{isomorphism} in \(\Sh_{X}(\category{C})\).
          Then there exists \(\psi: \sheaf{G} \to \sheaf{F}\) such that
          \begin{equation*}
            \psi \circ \phi = \id_{\sheaf{F}} \qquad \phi \circ \psi = \id_{\sheaf{G}}
          \end{equation*}
          Therefore since \hyperref[def:stalk]{taking stalks is functorial} for all \(p \in X\) we have
          \begin{equation*}
            \psi_{p} \circ \phi_{p} = \id_{\sheaf{F}_{p}} \qquad \phi_{p} \circ \psi_{p} = \id_{\sheaf{G}_{p}}
          \end{equation*}
          i.e: \(\phi_{p}: \sheaf{F}_{p} \to \sheaf{G}_{p}\) is an \hyperref[def:categorical-isomorphism]{isomorphism} with inverse \(\psi_{p}\).
    \item[\(\impliedby\)] Note that we only need to check the case where \(\category{C} = \Set\) since the \hyperref[def:functor]{functor} \(F\colon \category{C} \to \Set\) reflects \hyperref[def:categorical-isomorphism]{isomorphisms}.
          To do this, we will see that for all open subsets \(U \subseteq X\) the map \(\phi(U)\colon \sheaf{F}(U) \to \sheaf{G}(U)\) is injective and surjective.

          To see injectivity, let \(s, t \in \sheaf{F}(U)\) such that \(\phi(U)(s) = \phi(U)(t)\).
          For all \(p \in U\) we have
          \begin{equation*}
            \phi_{p}(s_{p}) = \left( \phi(U)(s) \right)_{p} = \left( \phi(U)(t) \right)_{p} = \phi_{p}(t_{p})
          \end{equation*}
          Therefore since \(\phi_{p}\) is injective \(s_{p} = t_{p}\).
          Using \hyperref[lemma:sheaf-determined-by-germs]{sheaves are determined by its germs} we get \(s = t\) and thus \(\phi(U)\) is injective.

          To see surjectivity consider \(t \in \sheaf{G}(U)\).
          For all \(p \in U\) since \(\phi_{p}\) is surjective there exists \(\sigma_{p} \in \sheaf{F}_{p}\) such that \(\phi_{p}(\sigma_{p}) = t_{p}\).
          By \hyperref[def:stalk]{definition of stalk} there exists an open neighbourhood \(V_{p} \subseteq U\) of \(p\) and a section \(s^{(p)} \in \sheaf{F}(V_{p})\) such that \(s_{p}^{(p)} = \sigma_{p}\).

          By shrinking \(V_{p}\) if necessary we can assume that
          \begin{equation*}
            \phi(V_{p})(s^{(p)}) = t|_{V_{p}}
          \end{equation*}
          Now consider two such neighbourhoods \(V_{p}\) and \(V_{q}\).
          On the intersection we have
          \begin{equation*}
            \phi(V_{p} \cap V_{q})(s^{(p)}|_{V_{p} \cap V_{q}}) = t|_{V_{p} \cap V_{q}} = \phi(V_{p} \cap V_{q})(s^{(q)}|_{V_{p} \cap V_{q}})
          \end{equation*}
          Therefore since \(\phi(V_{p} \cap V_{q})\) is injective \(s^{(p)}|_{V_{p} \cap V_{q}} = s^{(p)}|_{V_{p} \cap V_{q}}\).
          By the \hyperref[axiom:sheaf-gluability]{gluability axiom} there exists a unique section \(s \in \sheaf{F}(U)\) such that \(s|_{V_{p}} = s^{(p)}\) for all \(p \in U\).
          Finally since \(\phi(U)(s)|_{V_{p}} = \phi(V_{p})(s^{(p)}) = t|_{V_{p}}\) for all \(p \in U\) the uniqueness of the \hyperref[axiom:sheaf-gluability]{gluability axiom} now on \(\sheaf{G}\) implies that \(\phi(U)(s) = t\) and thus \(\phi(U)\) is surjective. \qedhere
  \end{itemize}
\end{proof}



\section{Sheafification}

We have established the definition of \hyperref[def:presheaf]{presheaves} and \hyperref[def:sheaf]{sheaves}. It is clear from the definitions that every \hyperref[def:sheaf]{sheaf} is a \hyperref[def:presheaf]{presheaf}. This defines a forgetful \hyperref[def:functor]{functor}:
\begin{equation*}
    \iota\colon \Sh_{X}(\category{C}) \hookrightarrow \PSh_{X}(\category{C})
\end{equation*}
We want to associate to each \hyperref[def:presheaf]{presheaf} a \hyperref[def:sheaf]{sheaf} in a way that \textit{respects} this forgetful functor.

In the language of category-theory, this is expressed by an \hyperref[def:adjoint-functors]{adjunction}.
\begin{definition}[Sheafification] \label{def:sheafification}
    Let \(X\) be a topological space and \(\category{C}\) a \hyperref[def:category]{category}.
    The \textbf{sheafification} (or associated sheaf) \hyperref[def:functor]{functor}, denoted by \(-^{\#}\colon \PSh_{X}(\category{C}) \to \Sh_{X}(\category{C})\), is the \hyperref[def:adjoint-functors]{left adjoint} to the inclusion functor \(\iota\colon \Sh_{X}(\category{C}) \hookrightarrow \PSh_{X}(\category{C})\).

    That is, for every \hyperref[def:presheaf]{presheaf} \(\mathcal{F} \in \PSh_{X}(\category{C})\) and every \hyperref[def:sheaf]{sheaf} \(\sheaf{G} \in \Sh_{X}(\category{C})\), there is a natural bijection:
    \begin{equation*}
        \Hom_{\Sh_{X}(\category{C})}(\mathcal{F}^{\#}, \sheaf{G}) \cong \Hom_{\PSh_{X}(\category{C})}(\mathcal{F}, \iota\sheaf{G})
    \end{equation*}

    The unit of this adjunction gives a canonical morphism \(\theta\colon \mathcal{F} \mapsto \mathcal{F}^{\#}\) in \(\PSh_{X}(\category{C})\) which verifies that for any \hyperref[def:presheaf]{presheaf morphism} \(\phi\colon \mathcal{F} \to \sheaf{G}\) where \(\mathcal{F}\) is a \hyperref[def:presheaf]{presheaf} and \(\sheaf{G}\) is a \hyperref[def:sheaf]{sheaf} there exists a unique morphism \(\psi\colon \mathcal{F}^{\#} \to \sheaf{G}\) such that \(\psi \circ \theta = \phi\).
\end{definition}

For this notion to be useful, we must prove that it exists.
We will do this by constructing the \hyperref[def:sheafification]{sheafification} of a \hyperref[def:presheaf]{presheaf} \(\mathcal{F}\), but before doing so we fix some terminology on constructing \hyperref[def:sheaf]{sheaves} from a \textit{bigger} one.

\begin{definition}\label{def:types-of-morphisms-pre-sheaves}
  Let \(X\) be a topological space.
  A \hyperref[def:presheaf]{presheaf} \(\mathcal{F}\) is a \textbf{subpresheaf} of a \hyperref[def:presheaf]{presheaf} \(\mathcal{G}\) if for all open sets \(U \subseteq X\) we have \(\mathcal{F}(U) \subseteq \mathcal{G}(U)\) and the restriction maps of \(\mathcal{G}\) induce the restriction maps of \(\mathcal{F}\).
  If \(\sheaf{F}\) and \(\sheaf{G}\) are \hyperref[def:sheaf]{sheaves} we say that \(\sheaf{F}\) is a \textbf{subsheaf} of \(\sheaf{G}\).
    % \item A \hyperref[def:presheaf]{morphism of presheaves} \(\phi\colon \mathcal{F} \to \mathcal{G}\) is \textbf{injective} if for all open sets \(U \subseteq X\) the map \(\mathcal{F}(U) \to \mathcal{G}(U)\) is injective.
    % \item A \hyperref[def:presheaf]{morphism of presheaves} \(\phi\colon \mathcal{F} \to \mathcal{G}\) is \textbf{surjective} if for all open sets \(U \subseteq X\) the map \(\mathcal{F}(U) \to \mathcal{G}(U)\) is surjective.
    % \item A \hyperref[def:sheaf]{morphism of sheaves} \(\varphi\colon \sheaf{F} \to \sheaf{G}\) is \textbf{injective} if for all open sets \(U \subseteq X\) the map \(\sheaf{F}(U) \to \sheaf{G}(U)\) is injective.
    % \item A \hyperref[def:sheaf]{morphism of sheaves} \(\varphi\colon \sheaf{F} \to \sheaf{G}\) is \textbf{surjective} if for all open sets \(U \subseteq X\) and every section \(s \in \sheaf{G}(U)\) there exists an open covering \(U = \bigcup_{i \in \mathcal{I}} U_{i}\) such that \(s|_{U_{i}}\) is in the image of \(\sheaf{F}(U_{i}) \to \sheaf{G}(U_{i})\) for all \(i\).
\end{definition}


\begin{construction}[Godement construction of sheafification] \label{const:sheafification-godement}
  Let \(\mathcal{F} \in \PSh_{X}(\category{\category{C}})\).

  For every open set \(U \subseteq X\), let
  \begin{equation*}
    \Lambda(\mathcal{F})(U) \coloneq \prod_{p \in U} \mathcal{F}_{p}
  \end{equation*}
  Considering for \(V \subseteq U\) the morphism
  \begin{align*}
    \res_{U, V}\colon \Lambda(\mathcal{F})(U) &\to \Lambda(\mathcal{F})(V) \\
    \prod_{p \in U} \mathcal{F}_{p} &\mapsto \prod_{p \in V} \mathcal{F}_{p}
  \end{align*}
  as the restriction maps, we see that \(\Lambda(\mathcal{F})\) is a \(\category{\category{C}}\)-valued \hyperref[def:presheaf]{presheaf}.

  Let's check that \(\Lambda(\mathcal{F})\) verifies the \hyperref[def:sheaf]{sheaf} axioms.
  To do this, consider an open set \(U \subseteq X\).
  \begin{enumerate}[label=SH.\arabic*), ref=SH.\arabic*]
    \item Let \(s, t \in \mathcal{F}(U)\) be two sections such that there exists an open covering \(\{U_{i}\}_{i \in \mathcal{I}}\) such that \(\forall i \in \mathcal{I} \quad s|_{U_{i}} = t|_{U_{i}}\).

          Note that for any \(x \in U\) since there exists \(j \in \mathcal{I}\) such that \(x \in U_{j}\) we have that \(s_{x} = t_{x}\).
          Therefore, since this holds for all \(x \in U\), we have that \(s = t\).
    \item Let \(\{U_{i}\}_{i \in \mathcal{I}}\) be an open covering of \(U\) and let \(\{s_{i} \mid s_{i} \in \Lambda(\mathcal{F})(U_{i})\}_{i \in \mathcal{I}}\) be a family of sections such that for any \(i, j \in \mathcal{I}\) we have \(s_{i}|_{U_{i} \cap U_{j}} = s_{j}|_{U_{i} \cap U_{j}}\).

          Define a section \(s \in \Lambda(\mathcal{F})(U)\) by choosing, for each \(x \in U\), an index \(k \in \mathcal{I}\) such that \(x \in U_{k}\) and setting \(s_{x} = (s_{k})_{x}\).
          Note that this is well-defined by the compatibility condition of the family.
          Furthermore, for all \(i \in \mathcal{I}\) by construction we have that \(s|_{U_{i}} = s_{i}\).
  \end{enumerate}

  Note that there is a \hyperref[def:presheaf]{presheaf morphism} \(\theta\colon \mathcal{F} \to \Lambda(\mathcal{F})\) defined for any \(s \in \mathcal{F}(U)\) by sending \(s\) to its germ at every point:
  \begin{align*}
    \theta(U)(s) = (s_{p})_{p \in U} \in \prod_{p \in U} \mathcal{F}_{p}
  \end{align*}
  where \(s_{p}\) denotes the image of \(s\) in the stalk \(\mathcal{F}_{p}\).

  We define the \hyperref[def:sheafification]{sheafification} \(\mathcal{F}^{\#}\) of \(\mathcal{F}\) as the \hyperref[def:types-of-morphisms-pre-sheaves]{subpresheaf} of \(\Lambda(\mathcal{F})\) consisting of sections that are \textit{locally compatible} with \(\mathcal{F}\).
  Explicitly, for every open set \(U \subseteq X\) an element \(\sigma = (\sigma_{p})_{p \in U} \in \Lambda(\mathcal{F})(U)\) belongs to \(\mathcal{F}^{\#}(U)\) if and only if:
  \begin{quote}
    For every \(p \in U\), there exists an open neighborhood \(V \subseteq U\) of \(p\) and a section \(t \in \mathcal{F}(V)\) such that for all \(q \in V\), \(\sigma_{q} = t_{q}\) (as elements of \(\mathcal{F}_{q}\)).
  \end{quote}

  We need to check that our definition of \(\mathcal{F}^{\#}\) verifies the universal property of \hyperref[def:sheafification]{sheafification}.
  To do this, consider a \hyperref[def:set-valued-sheaf]{sheaf} \(\sheaf{G}\) and a \hyperref[def:presheaf]{morphism of presheaves} \(\phi\colon \mathcal{F} \to \sheaf{G}\).
  We wish to show that there exists a unique \hyperref[def:sheaf]{sheaf} morphism \(\psi\colon \mathcal{F}^{\#} \to \sheaf{G}\) such that \(\phi \circ \theta = \psi\).
  Diagramatically
  \begin{center}
    \begin{tikzcd}
      \mathcal{F} \arrow[d, "\theta"'] \arrow[r, "\phi"] & \sheaf{G} \\
      \mathcal{F}^{\#} \arrow[ur, dotted, "\psi"']
    \end{tikzcd}
  \end{center}

  Let's check existence first by defining \(\psi\) explicitly.
  Let \(U \subseteq X\) be an open set and \(s \in \mathcal{F}^{\#}(U)\).
  By definition there exists an open cover \(\{U_{i}\}_{i \in \mathcal{I}}\) of \(U\) and sections \(t_{i} \in \mathcal{F}(V_{i})\) such that \(s|_{U_{i}} = \theta(t_{i})\) for all \(i \in \mathcal{I}\).

  We want \(\psi\) to verify \(\psi(s)|_{U_{i}} = \psi(s|_{U_{i}}) = \psi(\theta_{t_{i}})\). This forces us to define
  \begin{equation*}
    \psi(\theta(t_{i})) \coloneq \phi(t_{i}) \in \sheaf{G}(U_{i})
  \end{equation*}

  Let's see that \(\{\phi(t_{i})\}_{i \in \mathcal{I}}\) can be glued in \(\sheaf{G}\).
  Note that for all \(y \in U_{i} \cap U_{j}\) we have \((t_{i})_{y} = (t_{j})_{y}\).
  Therefore using the \hyperref[def:stalk]{induced stalk morphism} \(\phi_{y}\colon \mathcal{F}_{y} \to \sheaf{G}_{y}\) we have
  \begin{equation*}
    \phi_{y}((t_{i})_{y}) = \phi_{y}((t_{j})_{y})
  \end{equation*}
  Therefore \((\phi(t_{i}))_{y} = (\phi(t_{j}))_{y}\) for all \(y \in U_{i} \cap U_{j}\), meaning that
  \begin{equation*}
    \phi(t_{i})|_{U_{i} \cap U_{j}} = \phi(t_{j})|_{U_{i} \cap U_{j}}
  \end{equation*}

  Applying \ref{axiom:sheaf-gluability} to this family we get that there is a unique section \(t \in \sheaf{G}(U)\) such that \(t|_{U_{i}} = \phi(t_{i})\) for all \(i \in \mathcal{I}\).
  Therefore we can define \(\psi_{U}(s) \coloneq t\), which verifies the universal property by construction.

  To check uniqueness, since \hyperref[thm:sheaf-morphism-determined-by-stalks]{sheaf morphisms are determined by their induced maps on stalks} and for all points \(x \in X\) we have \(\theta_{x} = \id\), \(\psi\) verifies \(\psi_{x} = \phi_{x}\) for all \(x \in X\), making it unique.
\end{construction}

\begin{remark}
  Another construction of \hyperref[def:sheafification]{sheafification} is the \textit{espace étalé} of \(\mathcal{F}\).
  Roughly speaking, the section of \(\mathcal{F}^{\#}\) on an open set \(U \subseteq X\) correspond to
  \begin{equation*}
    \mathcal{F}^{\#}(U) = \left\{ s\colon U \to \coprod_{p \in P} \mathcal{F}_{p} \text{ continuous } \mid \pi \circ s = \id_{U} \right\}
  \end{equation*}
  where \(\pi\colon \coprod_{p \in P} \mathcal{F}_{p} \to X\) is the obvious projection sending a germ \(s_{p} \in \mathcal{F}_{p}\) to its base point \(p \in X\) and \(\coprod_{p \in P} \mathcal{F}_{p}\) has the topology generated by \(\{s_{p} \in \mathcal{F}_{p} \mid p \in U\}_{U \in \Open(X)}\).

  More details about this construction can be found at \cite{vakilRISINGSEAFoundations2025}.
\end{remark}

The following properties are direct.

\begin{lemma}\label{lemma:sheafification-of-sheaf-is-sheaf}
  Let \(\sheaf{F}\) be a \hyperref[def:sheaf]{sheaf}.
  Then \((\iota\sheaf{F})^{\#} \cong \sheaf{F}\), i.e: \(-^{\#} \circ \iota \cong \id_{\Sh_{X}(\category{C})}\).
\end{lemma}
\begin{proof}
  Consider the identity map \(\id: \iota\sheaf{F} \to \sheaf{F}\).
  Since \(\iota\sheaf{F}\) is a \hyperref[def:presheaf]{presheaf} and \(\sheaf{F}\) is a \hyperref[def:sheaf]{sheaf} by \hyperref[def:sheafification]{definition of sheafification} there exists a unique morphism \(\psi: (\iota\sheaf{F})^{\#} \to \sheaf{F}\) such that \(\psi \circ \theta = \id\).
  Note that
  \begin{equation*}
    (\theta \circ \psi) \circ \theta = \theta \circ (\psi \circ \theta) = \theta
  \end{equation*}
  Applying now the \hyperref[def:sheafification]{definition of sheafification} to \(\theta: \sheaf{F} \to (\iota\sheaf{F})^{\#}\) we get a unique map \(h: (\iota\sheaf{F})^{\#} \to (\iota\sheaf{F})^{\#}\) such that \(h \circ \theta = \theta\).
  Therefore since \(\id \circ \theta = \theta\) and \((\theta \circ \psi) \circ \theta = \theta\) we have that \(\theta \circ \psi = \id\), which means that \(\theta\) is an \hyperref[def:categorical-isomorphism]{isomorphism}.
\end{proof}


\begin{restatable}{lemma}{stalksheafification}\label{lemma:stalk-commute-sheafification}
  Let \(X\) be a topological space and \(\mathcal{F}\) a \hyperref[def:presheaf]{\(\category{C}\)-valued presheaf} over \(X\).
  Then for all \(p \in X\) there is bijection of \hyperref[def:stalk]{stalks}
  \begin{equation*}
    \mathcal{F}_{p} \cong (\mathcal{F}^{\#})_{p}
  \end{equation*}
\end{restatable}
\begin{proof}
  Direct from the \hyperref[const:sheafification-godement]{construction of sheafification}.
\end{proof}

% \begin{remark}
%   We will later give another proof of \autoref{lemma:stalk-commute-sheafification} that does not use the \hyperref[const:sheafification-godement]{constructions of sheafification}.
% \end{remark}


\section{Sheaves in an abelian category}

We want to apply the machinery we have built in \autoref{part:homological-algebra} to \hyperref[def:sheaf]{sheaves}.
In this section, \(X\) will be a fixed topological space and \(\category{C}\) will be an \hyperref[def:abelian-category]{abelian category} for which \hyperref[def:sheafification]{sheafification} exists.

Note that by \autoref{prop:functor-category-is-abelian} \(\PSh_{X}(\category{C})\) is an \hyperref[def:abelian-category]{abelian category}.
We are going to see that \(\Sh_{X}(\category{C})\) is an \hyperref[def:abelian-category]{abelian category}.

We will do this by the means of the following category-theoretic auxiliary lemma.
\begin{lemma}\label{lemma:abelian-category-from-adjunction}
  Let \(a\colon \category{A} \to \category{B}\) and \(b: \category{B} \to \category{A}\) be two \hyperref[def:functor]{functors}.
  Suppose that
  \begin{enumerate}
    \item \(\category{A}\) and \(\category{B}\) are \hyperref[def:additive-category]{additive categories}, \(a, b\) are \hyperref[def:additive-functor]{additive functors} and there is an \hyperref[def:adjoint-functors]{adjunction} \(b \dashv a\).
    \item \(\category{B}\) is an \hyperref[def:abelian-category]{abelian category} and \(b\) is \hyperref[def:exact-functor]{left-exact}.
    \item \(b \circ a \cong \id_{\category{A}}\).
  \end{enumerate}
  Then \(\category{A}\) is \hyperref[def:abelian-category]{abelian}.
\end{lemma}
\begin{proof}
  Since \(\category{B}\) is an \hyperref[def:abelian-category]{abelian} by \autoref{lemma:abelian-admits-limits-colimits} all finite \hyperref[def:categorical-limit]{limits} and \hyperref[def:categorical-colimit]{colimits} exist en \(\category{B}\).

  Since \(b\) is \hyperref[def:adjoint-functors]{left-adjoint} by \autoref{ex:adjoint-functors-and-exactness} it is \hyperref[def:exact-functor]{right-exact} and therefore, \hyperref[def:exact-functor]{exact}.
  Analogously we have that \(a\) is \hyperref[def:exact-functor]{left-exact}.

  Let \(\varphi\colon B_{1} \to B_{2}\) be a morphism in \(\category{B}\).
  Since \(\category{B}\) is \hyperref[def:abelian-category]{abelian} there exist \(\ker_{\category{B}} \varphi\) and \(\coker_{\category{B}} \varphi\).

  Since \hyperref[lemma:characterization-exact-functor]{exact functors preserve kernels and cokernels} we have that
  \begin{equation*}
    b(\ker_{\category{B}} \varphi) \cong \ker_{\category{A}}(b\varphi) \qquad b(\text{coker}_{\category{B}} \varphi) \cong \text{coker}_{\category{A}}(b\varphi)
  \end{equation*}
  
  Therefore we have that every morphism of the form \(b\varphi\) in \(\category{A}\) admits a \hyperref[def:categorical-kernel]{kernel} and a \hyperref[def:categorical-cokernel]{cokernel}.

  Note that, since \(b \circ a \cong \id_{\category{A}}\), every morphism \(\psi\) in \(\category{A}\) is of the form \(\psi = b\varphi\).
  Thus all \hyperref[def:categorical-kernel]{kernels} and \hyperref[def:categorical-cokernel]{cokernels} exist in \(\category{A}\) and
  \begin{align}\label{eq:ker-coker-lemma}
    \ker_{\category{A}} \psi &\cong b\ker_{\category{B}}(a\psi) \cong \ker_{\category{B}} b(a(\psi)) = \ker_{\category{b}} \psi \\
    \text{coker}_{\category{A}} \psi &\cong b\text{coker}_{\category{B}}(a\psi)
  \end{align}
  
  It only remains to show that \(\text{Coim}_{\category{A}} \psi \cong \text{Im}_{\category{A}} \psi\).
  Since \(\category{A}\) has all finite \hyperref[def:categorical-limit]{limits} and \hyperref[def:categorical-colimit]{colimits} \(a\) is \hyperref[def:exact-functor]{left exact}.

  By \autoref{lemma:characterization-exact-functor} we have that \(a\) transforms \hyperref[def:categorical-kernel]{kernels} into \hyperref[def:categorical-kernel]{kernels}.
  Therefore
  \begin{align*}
    \text{coim}_{\category{A}} \psi &= \text{coker}_{\category{A}} \left( \ker_{\category{A}} \psi \to A_{1} \right)      & (\text{definition}) \\
            &\cong b\text{coker}_{\category{B}} \left( a(\ker_{\category{A}} \psi \to A_{1}) \right)  & \eqref{eq:ker-coker-lemma} \\
            &\cong b\text{coker}_{\category{B}} \left( \ker_{\category{B}}(a\psi) \to a(A_{1}) \right)  & (a\text{ preserves \hyperref[def:categorical-kernel]{kernels}}) \\
            &= b\text{Coim}_{\category{B}} (a\psi)                                & (\text{definition}) \\
            &\cong b\text{Im}_{\category{B}} (a\psi)                                  & (\category{B} \text{ is \hyperref[def:abelian-category]{abelian}}) \\
            &= b(\ker_{\category{B}} \left( a(A_{2}) \to \text{coker}_{\category{B}}(a\psi) \right) ) & (\text{definition}) \\
            &\cong \ker_{\category{A}} \left( (b \circ a)(A_{2}) \to b(\text{coker}_{\category{B}}(a\psi)) \right) & (b\text{ preserves \hyperref[def:categorical-kernel]{kernels}}) \\
            &\cong \ker_{\category{A}} \left( A_{2} \to b(\text{coker}_{\category{B}}(a\psi)) \right)   & b \circ a \cong \id_{\category{A}}\\
            &\cong \ker_{\category{A}} \left( A_{2} \to \text{coker}_{\category{A}} \psi \right)      & \eqref{eq:ker-coker-lemma} \\
            &= \text{Im}_{\category{A}} \psi                                     & (\text{definition}) & \qedhere
  \end{align*}
\end{proof}

We need to check that we can apply this lemma in our case.

\begin{lemma}\label{lemma:sheaves-category-is-additive}
  \(\Sh_{X}(\category{C})\) is an \hyperref[def:additive-category]{additive category} and the forgetful \hyperref[def:functor]{functor} \(\iota\colon \Sh_{X}(\category{C}) \to \PSh_{X}(\category{C})\) is \hyperref[def:additive-functor]{additive}.
\end{lemma}
\begin{proof}
  Let's check first that \(\Sh_{X}(\category{C})\) satisfies \hyperref[def:additive-category]{additive category axioms}.
  Recall that \ref{axiom:additive-coproducts} is equivalent to \ref{axiom:additive-products} by \autoref{lemma:products-coproducts-additive-category}.

  We will use the \hyperref[rem:sheaves-as-equalisers]{characterization of sheaves in terms of equalisers}.
  \begin{enumerate}
    \item[\ref{axiom:additive-hom}] Since \(\Hom_{\Sh_{X}(\category{C})}(\sheaf{F}, \sheaf{G}) = \Hom_{\PSh_{X}(\category{C})}(\sheaf{F}, \sheaf{G})\) and \(\PSh_{X}(\category{C})\) is \hyperref[def:additive-category]{additive}, \(\Hom_{\Sh_{X}}\) admits the structure of an abelian group.
    \item[\ref{axiom:additive-zero}] Note that the \hyperref[def:zero-object]{zero object} in \(\PSh_{X}(\category{C})\) is given by the \hyperref[def:constant-presheaf]{constant presheaf} \(\underline{0}_{X}^{\pre}\).
          Note that the sequence
          \begin{center}
            \begin{tikzcd}
              \underline{0}_{X}^{\pre}(U) = 0  \arrow[r] & \displaystyle\prod_{i \in \mathcal{I}} \underline{0}_{X}^{\pre}(U_{i}) = 0 \arrow[r, shift left] \arrow[r, shift right] & \displaystyle\prod_{(i, j) \in \mathcal{I} \times \mathcal{I}} \underline{0}_{X}^{\pre}(U_{i} \cap U_{j}) = 0
            \end{tikzcd}
          \end{center}
          is trivially an \hyperref[def:equaliser]{equaliser} for all open covers \(U = \bigcup_{i \in \mathcal{I}} U_{i}\).
          Therefore \(\underline{0}_{X}^{\pre}\) is a \hyperref[def:sheaf]{sheaf} and thus \(\Sh_{X}(\category{C})\) admits a \hyperref[def:zero-object]{zero object}.
    \item[\ref{axiom:additive-products}] Let \(\sheaf{F}, \sheaf{G} \in \Sh_{X}(\category{C})\).
          We define their \hyperref[def:categorical-product]{product} elementwise as
          \begin{equation*}
            \left( \sheaf{F} \times \sheaf{G} \right)(U) = \sheaf{F}(U) \times \sheaf{G}(U)
          \end{equation*}
          Since this coincides with the \hyperref[def:categorical-product]{product} in \(\PSh_{X}(\category{C})\) it is a \hyperref[def:presheaf]{presheaf}.
          Furthermore for any open cover \(U = \bigcup_{i \in \mathcal{I}} U_{i}\) we have that
          \begin{align*}
            &(\sheaf{F} \times \sheaf{G})(U) = \sheaf{F}(U) \times \sheaf{G}(U) \to \prod_{i \in \mathcal{I}} (\sheaf{F} \times \sheaf{G})(U_{i}) = \prod_{i \in \mathcal{I}} \left(\sheaf{F}(U_{i})\right) \times \prod_{i \in \mathcal{I}} \left(\sheaf{G}(U_{i})\right) \\
            &\rightrightarrows \displaystyle\prod_{(i, j) \in \mathcal{I} \times \mathcal{I}} (\sheaf{F} \times \sheaf{G})(U_{i} \cap U_{j}) = \prod_{i \in \mathcal{I}} \left(\sheaf{F}(U_{i} \cap U_{j})\right) \times  \prod_{i \in \mathcal{I}} \left(\sheaf{G}(U_{i} \cap U_{j})\right)
          \end{align*}
          is an \hyperref[def:equaliser]{equaliser} since it can be seen as the ``product'' of two \hyperref[def:equaliser]{equalisers}.
          Therefore \(\sheaf{F} \times \sheaf{G}\) is a \hyperref[def:sheaf]{sheaf}.
  \end{enumerate}

  By the construction of the \hyperref[def:zero-object]{zero object} and \hyperref[def:categorical-product]{products} in \(\Sh_{X}(\category{C})\) we have that \(\iota\) preserves both.
  Therefore by \autoref{prop:additive-functor-characterization} \(\iota\) is \hyperref[def:additive-functor]{additive}.
\end{proof}

\begin{proposition}\label{prop:sheafification-is-exact}
  Let \(\category{C}\) be an \hyperref[def:abelian-category]{abelian category}.
  The \hyperref[def:sheafification]{sheafification functor} \(-^{\#}\colon \PSh_{X}(\category{C}) \to \Sh_{X}(\category{C})\) is \hyperref[def:exact-functor]{exact}.
\end{proposition}
\begin{proof}
  Since \(-^{\#}\) is \hyperref[def:adjoint-functors]{left adjoint} by \autoref{ex:adjoint-functors-and-exactness} it is \hyperref[def:exact-functor]{right exact}.

  To see that it is \hyperref[def:exact-functor]{left exact} we will use \autoref{lemma:characterization-exact-functor}.
  In particular we are going to see that \hyperref[def:sheafification]{sheafification} commutes with \hyperref[def:categorical-limit]{finite limits}.

  Let \(\mathcal{F}\colon \category{I} \to \PSh_{X}(\category{C})\) be a \hyperref[def:diagram]{diagram} in \(\PSh_{X}(\category{C})\) with \(\category{I}\) finite.

  We need to see that
  \begin{equation*}
    \left( \lim_{i \in \category{I}} \mathcal{F}_{i} \right)^{\#} \cong \lim_{i \in \category{I}} \mathcal{F}_{i}^{\#}
  \end{equation*}
  Note that since \(\iota\colon \Sh_{X}(\category{C}) \to \PSh_{X}(\category{C})\) is \hyperref[def:adjoint-functors]{right adjoint} it \hyperref[thm:adjoints-preserve-limits]{preserves limits}. Therefore \hyperref[def:categorical-limit]{limits} of \hyperref[def:sheaf]{sheaves} are computed exactly as \hyperref[def:categorical-limit]{limits} of \hyperref[def:presheaf]{presheaves}.

  By \hyperref[def:categorical-limit]{definition of limit} we have a morphism
  \begin{equation*}
    \phi\colon \left( \lim_{i \in \category{I}} \mathcal{F}_{i} \right)^{\#} \to \lim_{i \in \category{I}} \mathcal{F}_{i}^{\#}
  \end{equation*}
  To see that this is an \hyperref[def:categorical-isomorphism]{isomorphism} we will see that the \hyperref[def:stalk-map]{induced stalk maps} \(\phi_{p}\) are \hyperref[def:categorical-isomorphism]{isomorphisms} for all \(p \in X\), i.e: we want to show that
  \begin{equation*}
    \left( \left( \lim_{i \in \category{I}} \mathcal{F}_{i} \right)^{\#} \right)_{p} \cong \left( \lim_{i \in \category{I}} \mathcal{F}_{i}^{\#} \right)_{p}
  \end{equation*}


  On the one hand since \hyperref[lemma:stalk-commute-sheafification]{stalks commute with sheafification} we have
  \begin{equation*}
    \left( \left( \lim_{i \in \category{I}} \mathcal{F}_{i} \right)^{\#} \right)_{p} \cong \left( \lim_{i \in \category{I}} \mathcal{F}_{i} \right)_{p}
  \end{equation*}
  Furthermore since \hyperref[rem:stalk-commute-finite-limits]{stalks commute with finite limits}
  \begin{equation*}
    \left( \lim_{i \in \category{I}} \mathcal{F}_{i} \right)_{p} \cong \lim_{i \in \category{I}} \left(\mathcal{F}_{i}\right)_{p}
  \end{equation*}

  On the other hand since \hyperref[rem:stalk-commute-finite-limits]{stalks commute with finite limits}
  \begin{equation*}
    \left(\lim_{i \in \category{I}} \mathcal{F}_{i}^{\#} \right)_{p} \cong \lim_{i \in \category{I}} \left(\mathcal{F}_{i}^{\#}\right)_{p}
  \end{equation*}
  Furthermore since \hyperref[lemma:stalk-commute-sheafification]{stalks commute with sheafification} we have
  \begin{equation*}
    \lim_{i \in \category{I}} \left(\left( \mathcal{F}_{i} \right)^{\#}\right)_{p} \cong \lim_{i \in \category{I}} \left(\mathcal{F}_{i}\right)_{p}
  \end{equation*}
  Therefore
  \begin{equation*}
    \left( \left( \lim_{i \in \category{I}} \mathcal{F}_{i} \right)^{\#} \right)_{p} \cong \lim_{i \in \category{I}} \left(\mathcal{F}_{i}\right)_{p} \cong \left(\lim_{i \in \category{I}} \mathcal{F}_{i}^{\#} \right)_{p} \qedhere
  \end{equation*}
\end{proof}

We finally have everything we need to prove that \(\Sh_{X}(\category{C})\) is \hyperref[def:abelian-category]{abelian}.

\begin{theorem}\label{thm:sheaf-category-is-abelian}
  If \(\category{C}\) is \hyperref[def:abelian-category]{abelian} then the \hyperref[def:sheaf]{category of \(\category{C}\)-valued sheaves} \(\Sh_{X}(\category{C})\) over a topological space \(X\) is \hyperref[def:abelian-category]{abelian}.
  Furthermore if \(\varphi\colon \sheaf{F} \to \sheaf{G}\) is a morphism in \(\Sh_{X}(\category{C})\) its \hyperref[def:categorical-kernel]{kernel} is the same as its \hyperref[def:categorical-kernel]{kernel} as a map of \hyperref[def:presheaf]{presheaves} \(\ker \iota\varphi\) and its \hyperref[def:categorical-cokernel]{cokernel} is the \hyperref[def:sheafification]{sheafification} of its \hyperref[def:categorical-cokernel]{cokernel} as a map of \hyperref[def:presheaf]{presheaves}:
  \begin{align*}
    \ker \varphi = \ker \iota\varphi \qquad \coker \varphi = \left( \coker \iota\varphi \right)^{\#}
  \end{align*}
  and if \(p \in X\) then
  \begin{equation*}
    \left( \ker \varphi \right)_{p} = \ker (\varphi_{p}) \qquad \left( \coker \varphi \right)_{p} = \coker (\varphi_{p})
  \end{equation*}
\end{theorem}
\begin{proof}
  Direct from \autoref{lemma:abelian-category-from-adjunction} with
  \begin{align*}
    \category{A} \coloneq \Sh_{X}(\category{C}) &\quad \category{B} \coloneq \PSh_{X}(\category{B}) \\
    a \coloneq \iota\colon \Sh_{X}(\category{C}) \hookrightarrow \PSh_{X}(\category{C}) \text{ forgetful} &\quad b \coloneq -^{\#}\colon \PSh_{X}(\category{C}) \to \Sh_{X}(\category{C}) \text{ \hyperref[def:sheafification]{sheafification}}
  \end{align*}
  We can indeed apply \hyperref[lemma:abelian-category-from-adjunction]{this lemma} since
  \begin{enumerate}
    \item \(\Sh_{X}(\category{C})\) is \hyperref[def:additive-category]{additive} by \autoref{lemma:sheaves-category-is-additive} and \(\PSh_{X}(\category{C})\) is \hyperref[def:additive-category]{additive} by \autoref{lemma:functor-category-is-additive}, \(\iota\) is \hyperref[def:additive-functor]{additive} by \autoref{lemma:sheaves-category-is-additive} and \(-^{\#}\) is \hyperref[def:additive-functor]{additive} by \autoref{prop:additive-functors-and-adjunctions}, and the \hyperref[def:adjoint-functors]{adjunction} \(-^{\#} \dashv \iota\) is given by \hyperref[def:sheafification]{definition of sheafification}.
    \item \(\PSh_{X}(\category{C})\) is \hyperref[def:abelian-category]{abelian} by \autoref{prop:functor-category-is-abelian} and \(-^{\#}\) is \hyperref[def:exact-functor]{left exact} by \autoref{prop:sheafification-is-exact}.
    \item \(-^{\#} \circ \iota = \id_{\Sh_{X}(\category{C})}\) comes from \autoref{lemma:sheafification-of-sheaf-is-sheaf}.
  \end{enumerate}

  The commutativity of \hyperref[def:stalk]{stalks} comes form \autoref{rem:stalk-commute-finite-limits} and \autoref{lemma:stalk-commute-sheafification}.
\end{proof}

Furthermore combining \hyperref[thm:sheaf-category-is-abelian]{the above characterization of kernels and cokernels} and \autoref{prop:sheaf-isomorphism-determined-at-stalks} we have the following criterion to check \hyperref[def:exact-sequence]{exactness}.
\begin{corollary}\label{cor:exactness-at-stalks}
  Let \(\category{C}\) be an \hyperref[def:abelian-category]{abelian category}.
  A \hyperref[def:complex]{complex} of \hyperref[def:sheaf]{\(\category{C}\)-valued sheaves} over \(X\)
  \begin{equation*}
    0 \to \sheaf{F} \to \sheaf{G} \to \sheaf{H} \to 0
  \end{equation*}
  is \hyperref[def:exact-sequence]{exact} in \(\Sh_{X}(\category{C})\) if and only if for all \(p \in X\)
  \begin{equation*}
    0 \to \sheaf{F}_{p} \to \sheaf{G}_{p} \to \sheaf{H}_{p} \to 0
  \end{equation*}
   is \hyperref[def:exact-sequence]{exact} in \(\category{C}\).
\end{corollary}

\end{document}
